% COVID-19 material removed from the paper on 2026-11-18, kept for later use.

% ===== main.tex: Section 7 =====
%======================================================================
\section{Real data: Italy's COVID-19 restriction zones}\label{sec:italy-regions}
%======================================================================

Models that guide decisions are calibrated once and then used while the world
changes. The question this application asks is therefore not whether a simple
epidemic model is wrong, since every such model eventually is, but \emph{when}
a model used to set restrictions stops being adequate, \emph{where}, and
\emph{in which direction}, while the restrictions themselves are chosen from the
data that the model is checked against. The application has two parts.
Synthetic series, generated from the twin under a data-driven zone policy,
show that conditioning on the zones is what keeps the monitor valid
(Section~\ref{sec:italy-synth}). The real data then show what the alarms say about the twin: when
they fire, and whether cases ran above or below its predictions
(Section~\ref{sec:italy-real}).

\subsection{Zones, data and twin}\label{sec:italy}

From 6~November 2020, each of Italy's 21 regions and autonomous provinces was
placed every week, by government ordinance, in one of four zones: white (almost
no restrictions), yellow (curfew and evening closure of restaurants), orange
(restaurants closed and travel between municipalities restricted) or red (a
near-lockdown). The assignment followed published rules based on incidence,
the reproduction number and hospital indicators, so each action is an
explicit, discrete decision taken from the data, and the actions change often:
392 changes of colour, a median of 19 per region. Daily new positives by region
come from the Civil Protection Department \citep{dpc2020covid}, and the daily
colour of each region from a public compilation of the ordinances
\citep{imcatta2022restrizioni}.

The twin is a renewal model, the model behind the usual estimates of the
reproduction number \citep{cori2013new}. Region $r$'s daily count is negative
binomial with overdispersion $k$ and mean
\[
\mu_{rt} = \exp\big(a_r + \beta_{\mathrm{zone}(r,\,t-7)} + \varepsilon_{r,\mathrm{week}(t)}\big)\,
\delta_{\mathrm{wd}(t)} \sum_{s=1}^{21} w_s\, y_{r,t-s},
\]
where the weights $w_s$ are a gamma generation interval (mean 6.5, standard
deviation 4 days), $a_r$ is a regional baseline, $\beta$ are zone effects
relative to yellow, shared across regions, acting with a seven-day lag and
constrained so that red is at least as restrictive as orange and orange as
yellow, $\delta$ are day-of-week factors, and
$\varepsilon \sim \mathcal N(0, \sigma^2)$ is a shock common to the seven days
of a simulated week. The twin is calibrated on the first weeks of the zone
system, 13~November to 20~December 2020, by pooled maximum likelihood; $\sigma$
is chosen so that one-week-ahead 90\% intervals cover 90\% of region-weeks in
that window, which gives $\sigma = 0.10$. Orange multiplies the reproduction
number by 0.88 and red by 0.79; white zones began only in 2021, so the twin
treats them as yellow. From 21~December 2020 to 31~December 2021 each region is
monitored every week: its twin simulates the next seven days $M = 16$ times
from the real history and under the zone actually assigned, which is known when
the week starts because of the seven-day lag. After an alarm the region's
baseline is recalibrated on its last eight weeks and a new monitor starts.
Under the twin the regions are conditionally independent given their pasts, so
the average of the 21 regional wealth processes (frozen twin, no
recalibration) is a test martingale for the null that the twin holds in every
region, and e-BH applied to the regional e-values identifies the regions whose
twin fails. The twin deliberately omits variants, vaccination, seasonality,
testing volume and movement between regions: it is a plausible, standard model
of the kind used to project demand during the pandemic, and the question is
when and where it stops being adequate.

\subsection{Validity where the truth is known}\label{sec:italy-synth}

We first generate the 21 regional series from the twin, with zones assigned
every week by a documented rule that mimics the official one (red, orange,
yellow or white when the previous week's cases exceed 250, 150 or 50 per
100\,000 inhabitants, or none of these), with the estimated zone effects and
with stronger ones (orange $\times 0.85$, red $\times 0.70$), and a variant
that multiplies transmission by 1.3 in every region from week 8
(Table~\ref{tab:regions}). The zone-conditioned monitor is valid, with false
alarms in 1.5--1.9\% of region-runs and 1\% of national runs. A twin that
assumes every region is always yellow, and so ignores the actions, raises false
alarms in 21--23\% of region-runs, and the national monitor built on it rejects
a true null in every run. The plug-in bet and a $t$-test recomputed weekly
raise false alarms in about half and a third of region-runs. The national
monitor detects the variant in every run after a median of eight weeks.

\begin{table}[t]
\centering
\caption{Synthetic regional series with data-driven zones (53 weeks,
21 regions). Regional: share of region-runs with a false alarm. National: share
of runs in which the average-of-wealths monitor rejects. Variant (estimated
effects): detection by the national monitor and median delay in weeks.}
\label{tab:regions}
\begin{tabular}{lrrrrr}
\toprule
 & \multicolumn{2}{c}{estimated effects} & \multicolumn{2}{c}{stronger effects} & variant \\
monitor & regional & national & regional & national & national \\
\midrule
MC e-process, zone-conditioned twin & 0.019 & 0.01 & 0.015 & 0.01 & 1.00, 8 \\
twin assuming ``always yellow'' & 0.210 & 1.00 & 0.225 & 1.00 & 0.87, 8 \\
plug-in bet & 0.534 & -- & 0.508 & -- & -- \\
$t$-test, recomputed weekly & 0.319 & -- & 0.297 & -- & -- \\
\bottomrule
\end{tabular}
\end{table}

\subsection{Results on the real data}\label{sec:italy-real}

\begin{figure}[t]
\centering
\includegraphics[width=\textwidth]{exp9_italy_regions.png}
\caption{Italy's regional colour zones, December 2020 to December 2021. Top:
the zone assigned to each region (cells) and the monitor's alarms (triangles),
after each of which the region's twin is recalibrated. Bottom: the national
monitor, the average of the 21 regional wealths for the frozen twin, with the
threshold $\log 20$.}
\label{fig:regions}
\end{figure}

Figure~\ref{fig:regions} shows the analysis of the real data. The national
monitor rejects the twin on 14~March 2021, and by the end of 2021 e-BH rejects
it in all 21 regions. There are 96 regional alarms, a median of five per
region, and only Trento never alarms. The e-process tests the twin in both
directions, and the direction of each alarm is informative
(Table~\ref{tab:directions}): the alarms follow the waves of 2021, and within
each period almost all point the same way. From January to early April, 10 of
13 alarms fire because cases grew faster than the twin predicted \emph{given
the zone actually assigned}, for example in Bolzano on 24~January, already red
and at 580 cases per 100\,000 a week, and in Umbria, Campania, Lazio, Piedmont
and Liguria in February and March; the national alarm of 14~March belongs to
this period. The twin was calibrated before the Alpha variant spread, and
conditioning on the zone rules out relaxed restrictions as the explanation,
which makes this the clearest evidence in the application. All 27 alarms of
mid-April to early July fire because cases were \emph{below} the twin while
zones were being relaxed: the twin has neither vaccination nor seasonality, and
these alarms are correct but not warnings. Twelve regions alarm in the week
ending 25~July 2021 and seven more by 8~August, all in white zones and all
above the twin, which had just been recalibrated on the low levels of June,
during the Delta wave. From late August to October, 19 of 20 alarms are below
the twin as the summer wave receded faster than predicted, and in November and
December 16 of 17 are above it again. The monitor therefore identifies, with a
guarantee that holds although the zones were chosen from the data, two ways in
which the twin fails: the effect of a restriction learned in one wave does not
transfer to the next when the virus changes, and drivers outside the twin
dominate once restrictions stop varying, since 63 of the 96 alarms fire in
white zones, which the twin had to treat as yellow. A monitor that bets only on
counts above the twin, through a score that gives no credit to
overprediction, is equally valid and would keep the alarms of the Alpha, Delta
and autumn waves as early warnings. The results above use the Gaussian logistic learner; the other learners are
valid on the synthetic series, and on the real data the plug-in predictives
agree with it on the national alarm and on e-BH (Supplement, Section~S5). An analysis of national daily
deaths from spring 2020 under the national stringency index, in the
Supplement, gives an
early warning of the second wave on 22~September 2020, four weeks before deaths
surged.

\begin{table}[t]
\centering
\caption{The 96 regional alarms of 2021 by period. Above: observed cases above
the twin's simulations in the alarm week. Incidence: median weekly cases per
100\,000 in the alarm week.}
\label{tab:directions}
\begin{tabular}{lrrrrl}
\toprule
period & alarms & above & incidence & in white & likely driver \\
\midrule
January -- early April & 13 & 10 & 173 & 0 & Alpha variant, third wave \\
mid-April -- early July & 27 & 0 & 19 & 8 & vaccination, season \\
mid-July -- mid-August & 19 & 19 & 47 & 19 & Delta, reopening \\
late August -- October & 20 & 1 & 43 & 19 & end of the summer wave \\
November -- December & 17 & 16 & 103 & 17 & autumn Delta wave, Omicron \\
\bottomrule
\end{tabular}
\end{table}




% ===== main.tex: Data availability =====
\section*{Data availability}
The real-data applications use public data only. Italy's national and
regional COVID-19 series are from the Dipartimento della Protezione Civile
\citep{dpc2020covid}, the national stringency index is from the Oxford
COVID-19 Government Response Tracker \citep{hale2021global}, both under the
CC-BY-4.0 licence, and the regional restriction tiers are from a public
compilation of the government ordinances \citep{imcatta2022restrizioni}. The
processed files used in the paper, with the date of download and the exact
source files, are in the \texttt{data} folder of the code repository. All
other data are simulated by the code below.


% ===== supplement.tex: national deaths section =====
%======================================================================
\section{National COVID-19 deaths from spring 2020}\label{sec:italy-national}

\paragraph{Data and twin.} Daily deaths are the differences of the Civil
Protection cumulative count \citep{dpc2020covid}; one negative reporting
correction (24~June 2020) is set to zero. They are available from February
2020, and they depend much less on testing than case counts. The policy is
the national stringency index of the Oxford COVID-19 Government Response
Tracker \citep{hale2021global}, which covers the first lockdown of March 2020,
the reopening of May 2020 and the autumn measures.

The twin is the renewal model of Section~7.1 of the main text for
a single national series, without zones or weekly shocks: daily deaths are
negative binomial with mean $R_t\,\delta_{\mathrm{wd}(t)}\sum_{s=1}^{21} w_s y_{t-s}$
and $\log R_t = a + b\,(\mathrm{stringency}_{t-21}/100 - \tfrac12)$, with a
21-day lag because deaths follow infections by about three weeks, longer than
the seven-day block, so that the policy driving a week is known when the week
starts. Calibrated on 1~April to 26~May 2020, the twin has a baseline
reproduction number of 0.78 at the mean stringency of 88 and overdispersion
37.

The policy effect $b$ is poorly identified: estimated freely, it ranges from
$-0.4$ to $-3$ (the lower bound) depending on the 2020 window, as spontaneous
behaviour change and immunity are confounded with the measures. We therefore
fix $b = -1$, a 59\% reduction in $R$ between no restrictions and stringency
90, which is moderate compared with the lockdown effects estimated by
\citet{flaxman2020estimating}. We report $b \in \{-0.5, -1, -2\}$, a range
bracketed by the free estimates. The remaining parameters are fitted on
1~April to 26~May 2020, and monitoring runs from 27~May 2020 to 31~December
2021. After each alarm the twin is recalibrated on the last eight weeks.

\paragraph{Validity where the truth is known.} We generate series from the
calibrated twin over the same 83 weeks, with the policy set every week by a
tier rule on deaths: stringency 80, 60 or 40 when the previous week's deaths
exceed 1500, exceed 500, or neither. A second scenario adds a variant that
multiplies transmission by 1.4 from week 10.

\begin{table}[t]
\centering
\caption{Synthetic national series of deaths from the twin calibrated on
spring 2020, with a data-driven tier policy (83 weeks; 200 null and 100
variant replications).}
\label{tab:italy}
\begin{tabular}{lrrrr}
\toprule
 & null & \multicolumn{3}{c}{variant from week 10} \\
monitor & false alarms & alarm before change & detected & median delay (weeks) \\
\midrule
MC e-process, BLR score & 0.015 & 0.00 & 1.00 & 4 \\
plug-in bet & 0.035 & 0.01 & 0.99 & 3 \\
twin with policy held fixed & 1.000 & 1.00 & 0.00 & -- \\
$t$-test, recomputed weekly & 0.405 & 0.19 & 0.81 & 5 \\
\bottomrule
\end{tabular}
\end{table}

Table~\ref{tab:italy} repeats the regional lesson with a continuous policy.
The monitor keeps false alarms at 1.5\% and detects the variant in every run,
after a median of four weeks. The twin that holds the policy at its
calibration level rejects a true null in every run, and the recomputed
$t$-test raises false alarms in 41\% of runs.

\begin{figure}[t]
\centering
\includegraphics[width=\textwidth]{exp8b_italy_deaths.png}
\caption{Italy, May 2020 to December 2021. Top: observed daily deaths (black)
and the twin's 90\% predictive band for each coming week (blue); orange lines
mark alarms, after which the twin is recalibrated on the last eight weeks.
Bottom: the monitor's log wealth, restarted at each recalibration, with the
threshold $\log 20$.}
\label{fig:italy}
\end{figure}

\paragraph{Results on the real data.} Figure~\ref{fig:italy} shows ten alarms
between May 2020 and December 2021. Four fire because deaths were above the
twin and six because they were below it (average standardised deviation over
the last three weeks before the alarm):
\begin{itemize}
  \item \textbf{below}, 30~June 2020: deaths fell faster after the reopening
        than the spring twin predicted;
  \item \textbf{above}, 22~September and 3~November 2020: the start and the
        growth of the second wave;
  \item \textbf{below}, 15~December 2020 and 5~January 2021: after the peak of
        the second wave;
  \item \textbf{below}, then \textbf{above}, 2 and 23~March 2021: the end of
        the winter decline and the third wave;
  \item \textbf{below}, 18~May and 8~June 2021: the decline of late spring,
        with vaccination;
  \item \textbf{above}, 17~August 2021: the summer rise.
\end{itemize}

\paragraph{An early warning: September 2020.} After the first alarm the twin
is recalibrated on June 2020, when deaths were falling. With stringency almost
constant (63--68), it predicts a further decline. Weekly deaths instead reach
a minimum of about 40 in late August and then rise: 64, 69, 97 and 128 in the
four weeks from 31~August, each in the upper part of the twin's range. The log
wealth goes from $-0.6$ on 1~September to 0.85, 1.23 and 3.26 on 22~September,
when it crosses $\log 20$. Weekly deaths then reached 795 in the week of
19~October and 2568 in the week of 2~November. The alarm therefore precedes
the October surge in deaths by four weeks, and the national decrees of October
and the regional zones of 6~November by three to six weeks. Because the policy
was held at the value in force, the alarm says that transmission was higher
than the restrictions explain. Two qualifications apply: deaths lag infections,
so a twin of cases or admissions would probably alarm earlier; and the
comparison with the decrees is one of timing, not a counterfactual, since
decision makers also followed cases. The monitor is also robust to an outlier
in this segment: the 158 deaths reported on 15~August 2020 against 5--10 on
the neighbouring days, a batch of late reports, lower the wealth (from $-0.30$
to $-0.46$) instead of raising an alarm, because a bounded bet loses at most
its stake.

\paragraph{A correction, not a forecast: December 2020.} The alarm of
15~December is the opposite case. The twin recalibrated after the November
alarm expected continued growth, but deaths fell below its 90\% band in the
weeks from 7 to 21~December (for example 4279 against a lower limit of 4318 in
the week of 14~December). A plausible reading is that the November measures,
including the regional zones, which the national stringency index barely
records, reduced transmission more than the twin allowed. The alarm is
therefore not an anticipation of the national red days of Christmas 2020; it
says that the twin overpredicted.

Each recalibrated twin lasts a median of six weeks. The first alarm falls
between 16~June and 14~July 2020 for $b \in \{-0.5, -1, -2\}$, with 10 to 12
alarms over the period. An analysis of national ICU admissions from December
2020 to December 2022 with a seven-day lag leads to the same conclusions
(repository, Experiment~8).

\paragraph{Caveats.}
\begin{itemize}
  \item \textbf{An alarm says that the twin is wrong, not why.} Candidate
        reasons include variants, vaccination, seasonality, and changes in
        testing and reporting.
  \item \textbf{The twins' likelihoods happen to be tractable,} which keeps the
        examples transparent. The monitor uses only simulation, so the same
        analysis applies unchanged to agent-based or compartmental models with
        latent states.
  \item \textbf{Recalibration and restart form a sequence of tests of
        different twins.} Each alarm is valid for the twin active at the time,
        not a simultaneous statement about all periods.
\end{itemize}


% ===== supplement.tex: regional learners =====
\subsection{The regional analysis with each learner}

The real-data results above use the Gaussian logistic learner.
Table~\ref{tab:regions-learners} repeats the analysis with the four
predictive-first learners. All are valid on the synthetic series and detect
the variant nationally after 6--7 weeks, against 8 for the Gaussian learner. On
the real data the plug-in predictives agree with it on the national alarm (14
March 2021, or two weeks earlier) and on e-BH, which rejects the twin in all 21
regions; the resampled predictives are more cautious, with about half as many
regional alarms and a national alarm in April or June, and with the logistic
loss e-BH rejects no region at the end of 2021. The Bayesian-bootstrap forest,
refitted every five weeks from the fifth, is valid but slow, with a national
delay of 15 weeks: each region's forest learns from at most 53 weekly points.

\begin{table}[t]
\centering
\caption{The regional analysis with each learner. Synthetic (estimated
effects): regional and national false-alarm rates under the null, and median
national delay in weeks under the variant. Real data: number of regional
alarms, date of the national alarm, and number of regions rejected by e-BH at
the end of 2021.}
\label{tab:regions-learners}
\begin{tabular}{lrrrrrr}
\toprule
 & \multicolumn{3}{c}{synthetic} & \multicolumn{3}{c}{real data} \\
learner & regional & national & delay & alarms & national alarm & e-BH \\
\midrule
Gaussian BLR, logistic & 0.019 & 0.01 & 8 & 96 & 14 Mar 2021 & 21 \\
predictive, logistic, plug-in & 0.014 & 0.00 & 6 & 118 & 28 Feb 2021 & 21 \\
predictive, logistic, resampling & 0.013 & 0.02 & 6 & 52 & 4 Apr 2021 & 0 \\
predictive, simulation likelihood, plug-in & 0.003 & 0.01 & 6 & 86 & 14 Mar 2021 & 21 \\
predictive, simulation likelihood, resampling & 0.005 & 0.01 & 7 & 45 & 6 Jun 2021 & 20 \\
Bayesian-bootstrap forest & 0.002 & 0.00 & 15 & 32 & 1 Aug 2021 & 0 \\
\bottomrule
\end{tabular}
\end{table}

